% !TEX program = xelatex
\documentclass[11pt,a4paper]{article}
\usepackage[margin=2.2cm,top=2cm]{geometry}
\usepackage{fontspec}
\usepackage{polyglossia}
\setdefaultlanguage{turkish}
\usepackage{amsmath,amssymb}
\usepackage{graphicx}
\usepackage{booktabs}
\usepackage{array}
\usepackage{needspace}
\usepackage{enumitem}
\usepackage{multicol}
\usepackage[most]{tcolorbox}
\usepackage{fancyhdr}
\usepackage{titlesec}
\usepackage{xcolor}
\usepackage[hidelinks]{hyperref}

\definecolor{anamavi}{RGB}{0,45,98}
\definecolor{cozumyesil}{RGB}{0,110,50}
\definecolor{uyarikirmizi}{RGB}{180,20,20}

\titleformat{\section}{\large\bfseries\color{anamavi}}{\thesection.}{0.5em}{}
\titleformat{\subsection}{\normalsize\bfseries\color{anamavi}}{\thesubsection}{0.5em}{}

\pagestyle{fancy}\fancyhf{}
\lhead{\small MMAT105 Genel Matematik}
\rhead{\small Hafta 3 -- Ders Notu}
\cfoot{\small\thepage}
\renewcommand{\headrulewidth}{0.4pt}

\newtcolorbox{kural}[1][Kural]{colback=anamavi!5,colframe=anamavi,title=#1,fonttitle=\bfseries,breakable,left=4pt,right=4pt,top=2pt,bottom=2pt}
\newtcolorbox{dikkat}[1][Dikkat]{colback=uyarikirmizi!5,colframe=uyarikirmizi,title=#1,fonttitle=\bfseries,breakable,left=4pt,right=4pt,top=2pt,bottom=2pt}
\newcounter{ornek}
\newenvironment{ornek}{\needspace{5\baselineskip}\refstepcounter{ornek}\par\medskip\noindent\textbf{\color{anamavi}Örnek \theornek.}\ }{\par}
\newenvironment{cozum}{\begin{tcolorbox}[colback=cozumyesil!6,colframe=cozumyesil,boxrule=0.5pt,left=4pt,right=4pt,top=2pt,bottom=2pt,breakable]\textbf{\color{cozumyesil}Çözüm.}\ }{\end{tcolorbox}}
\newenvironment{cozumk}{\begin{tcolorbox}[colback=cozumyesil!6,colframe=cozumyesil,boxrule=0.5pt,left=4pt,right=4pt,top=2pt,bottom=2pt]\textbf{\color{cozumyesil}Çözüm.}\ }{\end{tcolorbox}}
\newcounter{alis}
\newenvironment{alistirma}{\refstepcounter{alis}\par\medskip\needspace{9\baselineskip}\noindent\textbf{\color{anamavi}Alıştırma \thealis.}\ }{\par}
\newcommand{\kendinizoncesi}{\needspace{12\baselineskip}}
\newcommand{\ab}[1]{\overline{#1}}
\setlength{\parindent}{0pt}
\setlength{\parskip}{3pt}

\begin{document}

% ---------------- Başlık ----------------
\begin{minipage}[c]{\textwidth}
{\LARGE\bfseries\color{anamavi} MMAT105 Genel Matematik}\\[4pt]
{\large Hafta 3 Ders Notu: Bölünebilme Kuralları,\\ EBOB ve EKOK}\\[6pt]
{\small 2026--2027 Güz}
\end{minipage}

\vspace{4pt}\noindent\textcolor{anamavi}{\rule{\textwidth}{1pt}}

\textbf{Bu notun kullanımı.} Her bölümde önce kısa bir kural kutusu, ardından çözümlü örnekler yer alır. Örnekleri çözümüne bakmadan önce kendiniz deneyin. Notun sonunda önce çözümlü alıştırmalar, ardından cevap anahtarlı ``Kendiniz Çözün'' alıştırmaları yer alır. Bu notta \(\ab{ab}\), \(\ab{abc}\) gibi gösterimler rakamlarla yazılmış sayılardır; ilk rakam \(0\) olamaz.

% =====================================================
\section{Temel Bölünebilme Kuralları}

\begin{kural}[Bölünebilme ve kalan]
\renewcommand{\arraystretch}{1.3}
\begin{tabular}{@{}>{\raggedright\arraybackslash}p{0.07\linewidth}>{\raggedright\arraybackslash}p{0.44\linewidth}>{\raggedright\arraybackslash}p{0.43\linewidth}@{}}
\toprule
Bölen & Sayı bölünür, eğer \dots & Kalan \\\midrule
\(2\) & birler basamağı çift (\(0,2,4,6,8\)) ise & birler basamağının \(2\)'ye bölümünden kalan\\
\(5\) & birler basamağı \(0\) veya \(5\) ise & birler basamağının \(5\)'e bölümünden kalan\\
\(10\) & birler basamağı \(0\) ise & birler basamağındaki rakam\\
\(3\) & rakamları toplamı \(3\)'ün katı ise & rakamları toplamının \(3\)'e bölümünden kalan\\
\(9\) & rakamları toplamı \(9\)'un katı ise & rakamları toplamının \(9\)'a bölümünden kalan\\
\(4\) & son iki basamağın oluşturduğu sayı \(4\)'ün katı ise & son iki basamağın \(4\)'e bölümünden kalan\\
\(8\) & son üç basamağın oluşturduğu sayı \(8\)'in katı ise & son üç basamağın \(8\)'e bölümünden kalan\\
\(11\) & birler basamağından başlayarak \(+,-,+,-,\dots\) işaretleriyle alınan toplam \(11\)'in katı ise (\(0\) dahil) & bu toplamın \(11\)'e bölümünden kalan (toplam negatifse \(11\)'in katı eklenir)\\\bottomrule
\end{tabular}
\end{kural}

\begin{dikkat}
\(9\) ile bölünen sayı \(3\) ile, \(8\) ile bölünen sayı \(4\) ile de bölünür; tersleri her zaman doğru değildir (\(12\), \(3\) ile bölünür ama \(9\) ile bölünmez). Bilinmeyen rakam soruda \(0\) değerini de denemeyi unutmayın. \(11\) kuralında işaretler \textbf{sağdan} başlar: \(\ab{abcd}\rightarrow d-c+b-a\).
\end{dikkat}

\begin{ornek}
\(\ab{a3b}\) üç basamaklı sayısı hem \(10\) ile hem \(3\) ile bölünebiliyorsa \(a\) hangi değerleri alabilir?
\end{ornek}
\begin{cozum}
\(10\) ile bölünme: \(b=0\). \ Rakamlar toplamı \(a+3\) sayısı \(3\)'ün katı olmalı; \(a\neq 0\) (ilk rakam).
\[a+3\in\{6,9,12\}\ \Rightarrow\ a\in\{3,6,9\}\qquad(330,\ 630,\ 930)\]
\end{cozum}

\begin{ornek}
\(7\,256\,814\) sayısının \(3\) ile ve \(9\) ile bölümünden kalanlar kaçtır?
\end{ornek}
\begin{cozum}
Rakamlar toplamı: \(7+2+5+6+8+1+4=33\). \quad \(33=3\cdot 11\) \(\Rightarrow\) \(3\) ile kalan \(0\); \quad \(33=9\cdot 3+6\) \(\Rightarrow\) \(9\) ile kalan \(6\).
\end{cozum}

\begin{ornek}
\(\ab{4a3b}\) dört basamaklı sayısı \(9\) ile bölünebiliyorsa \(a+b\) toplamı hangi değerleri alabilir?
\end{ornek}
\begin{cozum}
Rakamlar toplamı \(7+a+b\) sayısı \(9\)'un katı olmalı. \(0\le a+b\le 18\) olduğundan \(7+a+b\in\{9,18\}\):
\[a+b=2\quad\text{veya}\quad a+b=11\]
(\(7+a+b=27\) için \(a+b=20\) olur; iki rakamın toplamı en çok \(18\)'dir.)
\end{cozum}

\begin{ornek}
\(\ab{12a4}\) dört basamaklı sayısı \(8\) ile bölünebiliyorsa \(a\) hangi değerleri alabilir?
\end{ornek}
\begin{cozum}
Son üç basamak \(\ab{2a4}\) \(8\)'in katı olmalıdır. \(200\) ile \(299\) arasında \(4\) ile biten \(8\)'in katları: \(224=8\cdot 28\) ve \(264=8\cdot 33\) (\(204, 244, 284\) değil).
\[a\in\{2,6\}\qquad(1224=8\cdot 153,\ \ 1264=8\cdot 158)\]
\end{cozum}

\begin{ornek}
\(\ab{5a17}\) dört basamaklı sayısının \(11\) ile bölümünden kalan \(3\) ise \(a\) kaçtır?
\end{ornek}
\begin{cozum}
Sağdan: \(7-1+a-5=a+1\). \ \(a+1\) sayısının \(11\)'e bölümünden kalan \(3\) olmalı; \(1\le a+1\le 10\) olduğundan \(a+1=3\Rightarrow a=2\).\\
Kontrol: \(5217=11\cdot 474+3\).
\end{cozum}

% =====================================================
\section{Bileşik Sayılarla Bölünebilme}

\begin{kural}[Aralarında asal çarpanlara ayırma]
Bir sayının bileşik bir sayıyla bölünebilmesi için bölen, \textbf{aralarında asal} (\(1\)'den başka ortak pozitif böleni olmayan) çarpanlara ayrılır ve sayı her çarpanla ayrı ayrı bölünmelidir:
\[6=2\cdot 3,\ \ 12=3\cdot 4,\ \ 15=3\cdot 5,\ \ 18=2\cdot 9,\ \ 24=3\cdot 8,\ \ 36=4\cdot 9,\ \ 45=5\cdot 9,\ \ 72=8\cdot 9\]
\end{kural}

\begin{dikkat}
\(12=2\cdot 6\) veya \(24=4\cdot 6\) gibi ayrıştırmalar \textbf{kullanılamaz}: çarpanlar aralarında asal değildir. Örnek: \(18\) sayısı \(2\) ile ve \(6\) ile bölünür, fakat \(12\) ile bölünmez.
\end{dikkat}

\begin{ornek}
\(\ab{a5b}\) üç basamaklı sayısı \(12\) ile bölünebiliyorsa bu koşulu sağlayan kaç farklı sayı vardır?
\end{ornek}
\begin{cozum}
\(12=3\cdot 4\). \ \(4\): \(\ab{5b}\in\{52,56\}\Rightarrow b\in\{2,6\}\).\\
\(3\): \(a+5+b\) sayısı \(3\)'ün katı.\\
\(b=2\): \(a+7\) \(3\)'ün katı \(\Rightarrow a\in\{2,5,8\}\) \quad (\(252,\ 552,\ 852\))\\
\(b=6\): \(a+11\) \(3\)'ün katı \(\Rightarrow a\in\{1,4,7\}\) \quad (\(156,\ 456,\ 756\))\\
Toplam \(6\) farklı sayı.
\end{cozum}

\begin{ornek}
\(\ab{3a16}\) dört basamaklı sayısı \(24\) ile bölünebiliyorsa \(a\) hangi değerleri alabilir?
\end{ornek}
\begin{cozum}
\(24=3\cdot 8\).\\
\(8\): \(\ab{a16}\) \(8\)'in katı olmalı. \(016, 216, 416, 616, 816\) sayıları \(8\)'in katıdır; \(116, 316, 516, 716, 916\) değildir \(\Rightarrow\) \(a\) çift.\\
\(3\): \(3+a+1+6=10+a\) \(3\)'ün katı \(\Rightarrow a\in\{2,5,8\}\).\\
Ortak: \(a\in\{2,8\}\) \quad (\(3216=24\cdot 134\), \ \(3816=24\cdot 159\)).
\end{cozum}

% =====================================================
\section{EBOB (En Büyük Ortak Bölen)}

\begin{kural}[EBOB]
\(\text{EBOB}(a,b)\): pozitif tam sayılar \(a\) ve \(b\)'nin ortak pozitif bölenlerinin en büyüğü.
\begin{itemize}[leftmargin=*,itemsep=1pt]
\item \textbf{Asal çarpan yöntemi:} Yalnızca \textbf{ortak} asal çarpanlar, \textbf{en küçük} üsleriyle çarpılır.
\item \textbf{Ortak bölme yöntemi:} Sayılar, hepsini birden bölen asal sayılarla bölünür; bu bölenlerin çarpımı EBOB'dur.
\item \(\text{EBOB}(a,b)=1\) ise \(a\) ve \(b\) \textbf{aralarında asaldır}.
\end{itemize}
\end{kural}

\needspace{11\baselineskip}
\begin{ornek}
\(\text{EBOB}(90,120,150)\) kaçtır?
\end{ornek}
\begin{cozumk}\par
\begin{minipage}[t]{0.40\linewidth}
\begin{tabular}{rrr|l}
\(90\) & \(120\) & \(150\) & \(2\)\\
\(45\) & \(60\) & \(75\) & \(3\)\\
\(15\) & \(20\) & \(25\) & \(5\)\\
\(3\) & \(4\) & \(5\) & (ortak bölen yok)
\end{tabular}
\end{minipage}\hfill
\begin{minipage}[t]{0.56\linewidth}
\[\text{EBOB}=2\cdot 3\cdot 5=30\]
Asal çarpan yöntemiyle: \(90=2\cdot 3^2\cdot 5\), \(120=2^3\cdot 3\cdot 5\), \(150=2\cdot 3\cdot 5^2\); ortak çarpanların en küçük üsleri \(2\cdot 3\cdot 5=30\).
\end{minipage}
\end{cozumk}

\begin{ornek}
\(A=2^5\cdot 3^2\cdot 7\) ve \(B=2^3\cdot 3^4\cdot 5\) olduğuna göre \(\text{EBOB}(A,B)\) sayısının kaç tane pozitif tam böleni vardır?
\end{ornek}
\begin{cozum}
Ortak asal çarpanlar \(2\) ve \(3\): \ \(\text{EBOB}(A,B)=2^3\cdot 3^2=72\).\\
Pozitif bölen sayısı (2. hafta): \((3+1)(2+1)=12\).
\end{cozum}

% =====================================================
\section{EKOK (En Küçük Ortak Kat) ve EBOB\,\(\cdot\)\,EKOK}

\begin{kural}[EKOK ve temel bağıntı]
\(\text{EKOK}(a,b)\): pozitif tam sayılar \(a\) ve \(b\)'nin pozitif ortak katlarının en küçüğü.
\begin{itemize}[leftmargin=*,itemsep=1pt]
\item \textbf{Asal çarpan yöntemi:} \textbf{Bütün} asal çarpanlar, \textbf{en büyük} üsleriyle çarpılır.
\item \textbf{Bölme tablosu:} Sayılar, en az birini bölen asal sayılarla hepsi \(1\) olana kadar bölünür (bölünmeyen sayı aynen aşağı yazılır); bütün bölenlerin çarpımı EKOK'tur.
\item İki sayı için: \(\text{EBOB}(a,b)\cdot\text{EKOK}(a,b)=a\cdot b\). \ Aralarında asal ise \(\text{EKOK}(a,b)=a\cdot b\).
\item \(\text{EBOB}(a,b)\le\min(a,b)\le\max(a,b)\le\text{EKOK}(a,b)\).
\end{itemize}
\end{kural}

\begin{dikkat}
\(\text{EBOB}\cdot\text{EKOK}=\) çarpım bağıntısı yalnızca \textbf{iki} sayı için geçerlidir (\(2,4,8\) için \(2\cdot 8\neq 2\cdot 4\cdot 8\)). Aralarında asal sayıların asal olması gerekmez: \(8\) ve \(15\) aralarında asaldır.
\end{dikkat}

\needspace{12\baselineskip}
\begin{ornek}
\(\text{EKOK}(20,24,30)\) kaçtır?
\end{ornek}
\begin{cozumk}\par
\begin{minipage}[t]{0.40\linewidth}
\begin{tabular}{rrr|l}
\(20\) & \(24\) & \(30\) & \(2\)\\
\(10\) & \(12\) & \(15\) & \(2\)\\
\(5\) & \(6\) & \(15\) & \(2\)\\
\(5\) & \(3\) & \(15\) & \(3\)\\
\(5\) & \(1\) & \(5\) & \(5\)\\
\(1\) & \(1\) & \(1\) &
\end{tabular}
\end{minipage}\hfill
\begin{minipage}[t]{0.56\linewidth}
\[\text{EKOK}=2^3\cdot 3\cdot 5=120\]
Asal çarpan yöntemiyle: \(20=2^2\cdot 5\), \(24=2^3\cdot 3\), \(30=2\cdot 3\cdot 5\); en büyük üsler \(2^3\), \(3\), \(5\).
\end{minipage}
\end{cozumk}

\begin{ornek}
İki pozitif tam sayının EBOB'u \(8\), EKOK'u \(96\)'dır. Sayılardan biri \(32\) ise diğeri kaçtır?
\end{ornek}
\begin{cozum}
\[32\cdot b=8\cdot 96=768\ \Rightarrow\ b=24\]
Kontrol: \(32=2^5\), \(24=2^3\cdot 3\) \(\Rightarrow\) \(\text{EBOB}=2^3=8\), \(\text{EKOK}=2^5\cdot 3=96\).
\end{cozum}

\begin{ornek}
\(a\) ve \(b\) pozitif tam sayılardır. \(\text{EBOB}(a,b)=6\) ve \(a+b=48\) ise kaç farklı \((a,b)\) sıralı ikilisi vardır?
\end{ornek}
\begin{cozum}
\(a=6x\), \(b=6y\) yazılırsa \(x\) ile \(y\) aralarında asal pozitif tam sayılardır ve \(6(x+y)=48\Rightarrow x+y=8\).\\
\((x,y)\): \((1,7),(3,5),(5,3),(7,1)\) uygun; \((2,6),(4,4),(6,2)\) aralarında asal değil.\\
Cevap: \(4\) sıralı ikili: \((6,42),(18,30),(30,18),(42,6)\).
\end{cozum}

% =====================================================
\section{EBOB ve EKOK Problemleri}

\begin{kural}[Hangisi?]
\textbf{EBOB:} büyük bir şeyi eşit parçalara bölme (kesme, fayans döşeme, eşit aralıklı dikme); ``en büyük parça'', ``en az parça''.\\
\textbf{EKOK:} küçük şeyleri birleştirme (birlikte çalma, yeniden karşılaşma, küçük parçalardan kare oluşturma, hep aynı kalanı veren sayı).\\
\textbf{Ağaç dikme:} Kapalı çevrede ağaç sayısı \(=\) çevre \(\div\) aralık; iki ucuna da ağaç dikilen açık yolda \(=\) uzunluk \(\div\) aralık \(+1\).\\
\textbf{Kalanlı sorular:} \(a,b,c\) ile bölündüğünde hep \(r\) kalan sayılar: \(\text{EKOK}\cdot k+r\). \ Kalanlar bölenlerden hep \(t\) eksikse: \(\text{EKOK}\cdot k-t\).
\end{kural}

\begin{ornek}
Boyları \(150\) cm, \(210\) cm ve \(270\) cm olan üç kumaş, hiç artmayacak biçimde eşit boyda parçalara kesilecektir. En az kaç parça elde edilir?
\end{ornek}
\begin{cozum}
\(\text{EBOB}(150,210,270)=30\) cm. \quad Parça sayısı: \(\dfrac{150}{30}+\dfrac{210}{30}+\dfrac{270}{30}=5+7+9=21\).
\end{cozum}

\begin{ornek}
Kenarları \(24\) cm ve \(36\) cm olan eş dikdörtgen kartonlar, aynı yönde ve boşluksuz yan yana dizilerek bir kare oluşturulacaktır. En az kaç karton gerekir?
\end{ornek}
\begin{cozum}
Karenin kenarı hem \(24\)'ün hem \(36\)'nın katı olmalı ve en küçük olmalı: \(\text{EKOK}(24,36)=72\) cm.
\[\frac{72}{24}\cdot\frac{72}{36}=3\cdot 2=6\ \text{karton}\]
\end{cozum}

\begin{ornek}
Kenar uzunlukları \(42\) m, \(56\) m ve \(70\) m olan üçgen biçimli bir tarlanın çevresine, her köşeye birer ağaç gelecek ve ardışık ağaçlar arası uzaklık eşit olacak biçimde ağaç dikilecektir. En az kaç ağaç gerekir?
\end{ornek}
\begin{cozum}
Aralık: \(\text{EBOB}(42,56,70)=14\) m. \quad Çevre: \(42+56+70=168\) m.\\
Kapalı çevre: \(168\div 14=12\) ağaç. \quad (Kenar kenar: \(3+4+5=12\) aralık \(=12\) ağaç.)
\end{cozum}

\begin{ornek}
\(5\) ile bölündüğünde \(3\), \(6\) ile bölündüğünde \(4\), \(8\) ile bölündüğünde \(6\) kalanını veren en küçük pozitif tam sayı kaçtır?
\end{ornek}
\begin{cozum}
Kalanların her biri bölenden \(2\) eksiktir: \(x+2\) sayısı \(5\), \(6\) ve \(8\)'in ortak katıdır.
\[x+2=\text{EKOK}(5,6,8)=120\ \Rightarrow\ x=118\]
Kontrol: \(118=5\cdot 23+3=6\cdot 19+4=8\cdot 14+6\).
\end{cozum}

\begin{ornek}
Bir kutudaki cevizler \(6\)'şar, \(9\)'ar ve \(12\)'şer sayıldığında her seferinde \(2\) ceviz artmaktadır. Kutudaki ceviz sayısı \(50\) ile \(100\) arasında ise kaç ceviz vardır?
\end{ornek}
\begin{cozum}
Ceviz sayısı \(x\) ise \(x-2\) sayısı \(\text{EKOK}(6,9,12)=36\)'nın katıdır: \(x\in\{2,38,74,110,\dots\}\).\\
\(50<x<100\) aralığındaki tek değer \(x=74\).
\end{cozum}

% =====================================================
\section{Özet: Formül Kartı}
\begin{center}
\renewcommand{\arraystretch}{1.35}
\begin{tabular}{@{}ll@{}}
\toprule
Konu & Formül / Kural\\\midrule
\(2,\ 5,\ 10\) & birler basamağı\\
\(3,\ 9\) & rakamlar toplamı (kalan da rakamlar toplamından)\\
\(4,\ 8\) & son iki / son üç basamak\\
\(11\) & sağdan \(+,-,+,-,\dots\) toplamı\\
Bileşik bölen & aralarında asal çarpanlar: \(12=3\cdot 4\), \(36=4\cdot 9\), \(72=8\cdot 9\) \ (\(2\cdot 6\) değil)\\
EBOB & ortak asal çarpanlar, en küçük üs\\
EKOK & bütün asal çarpanlar, en büyük üs\\
İki sayı & \(\text{EBOB}(a,b)\cdot\text{EKOK}(a,b)=a\cdot b\); \ aralarında asal \(\Rightarrow\) EBOB \(=1\)\\
Kesme, fayans, ağaç & EBOB; \ kapalı çevrede ağaç \(=\) çevre \(\div\) aralık\\
Birlikte çalma & EKOK\\
Hep \(r\) kalan / hep \(t\) eksik & \(\text{EKOK}\cdot k+r\) \ / \ \(\text{EKOK}\cdot k-t\)\\\bottomrule
\end{tabular}
\end{center}

% =====================================================
\newpage
% Hafta 3 -- alıştırma gövdesi (ders notu ve ayrı alıştırma kağıdı tarafından \input edilir)
\section{Çözümlü Alıştırmalar}
Her soruyu önce kendiniz deneyin, sonra çözümü okuyun.

\begin{alistirma}
\(\ab{6a3}\) üç basamaklı sayısı \(11\) ile bölünebiliyorsa \(a\) kaçtır?
\end{alistirma}
\begin{cozum}
Birler basamağından başlayarak \(+,-,+\): \ \(3-a+6=9-a\). Bu sayı \(11\)'in katı olmalıdır.\\
\(0\le a\le 9\) olduğundan \(0\le 9-a\le 9\); bu aralıktaki tek \(11\) katı \(0\)'dır: \(9-a=0\Rightarrow a=9\) \ (\(693=11\cdot 63\)).
\end{cozum}

\begin{alistirma}
\(\ab{2a5b}\) dört basamaklı sayısı \(5\) ile ve \(9\) ile bölünebiliyorsa \(a\) en çok kaçtır?
\end{alistirma}
\begin{cozum}
\(5\) ile bölünme: \(b\in\{0,5\}\). \quad \(9\) ile bölünme: rakamlar toplamı \(2+a+5+b=7+a+b\) sayısı \(9\)'un katı.\\
\(b=0\): \(7+a\in\{9\}\) (\(7\le 7+a\le 16\)) \(\Rightarrow a=2\) \quad (\(2250\))\\
\(b=5\): \(12+a\in\{18\}\) (\(12\le 12+a\le 21\)) \(\Rightarrow a=6\) \quad (\(2655\))\\
\(a\) en çok \(6\)'dır. \ Kontrol: \(2655=5\cdot 531=9\cdot 295\).
\end{cozum}

\begin{alistirma}
\(\ab{3a4b}\) dört basamaklı sayısı \(36\) ile bölünebiliyorsa bu koşulu sağlayan kaç farklı sayı vardır?
\end{alistirma}
\begin{cozum}
\(36=4\cdot 9\) (aralarında asal çarpanlar; \(36=6\cdot 6\) veya \(3\cdot 12\) kullanılamaz).\\
\(4\) ile bölünme: \(\ab{4b}\in\{40,44,48\}\Rightarrow b\in\{0,4,8\}\).\\
\(9\) ile bölünme: \(3+a+4+b=7+a+b\) sayısı \(9\)'un katı.\\
\(b=0\): \(7+a=9\Rightarrow a=2\) \ (\(3240\)); \quad \(b=4\): \(11+a=18\Rightarrow a=7\) \ (\(3744\)); \quad \(b=8\): \(15+a=18\Rightarrow a=3\) \ (\(3348\)).\\
Cevap: \(3\) farklı sayı. \ (Kontrol: \(3240=36\cdot 90\), \(3348=36\cdot 93\), \(3744=36\cdot 104\).)
\end{cozum}

\begin{alistirma}
\(A=2^3\cdot 3^2\cdot 5\) ve \(B=2\cdot 3^3\cdot 5^2\) ise \(\text{EBOB}(A,B)\) ve \(\text{EKOK}(A,B)\) kaçtır?
\end{alistirma}
\begin{cozum}
EBOB: ortak asal çarpanlar en küçük üsleriyle: \ \(2^{1}\cdot 3^{2}\cdot 5^{1}=90\).\\
EKOK: bütün asal çarpanlar en büyük üsleriyle: \ \(2^{3}\cdot 3^{3}\cdot 5^{2}=8\cdot 27\cdot 25=5400\).\\
Kontrol: \(A=360\), \(B=1350\); \ \(\text{EBOB}\cdot\text{EKOK}=90\cdot 5400=486\,000=360\cdot 1350=A\cdot B\) \checkmark
\end{cozum}

\begin{alistirma}
\(a\) ve \(b\) aralarında asal iki pozitif tam sayıdır ve \(a\cdot b=90\)'dır. \(a+b\) en az kaçtır?
\end{alistirma}
\begin{cozum}
\(90=2\cdot 3^2\cdot 5\). \(a\) ile \(b\) aralarında asal olduğundan her asal çarpan bütün kuvvetiyle \textbf{yalnız birinde} bulunur; \(3^2=9\) bölünemez. Parçalar: \(2\), \(9\), \(5\).\\
Olası ikililer: \(\{1,90\}\to 91\), \ \(\{2,45\}\to 47\), \ \(\{5,18\}\to 23\), \ \(\{9,10\}\to 19\).\\
\(a+b\) en az \(9+10=19\). \ (\(\{6,15\}\) çarpımı \(90\) olsa da \(\text{EBOB}(6,15)=3\) olduğundan uygun değildir.)
\end{cozum}

\begin{alistirma}
Kenarları \(96\) m ve \(120\) m olan dikdörtgen bir bahçenin çevresine, her köşede birer ağaç olacak ve ardışık ağaçlar arası uzaklıklar eşit olacak biçimde ağaç dikilecektir. En az kaç ağaç gerekir?
\end{alistirma}
\begin{cozum}
En az ağaç için aralık en büyük olmalı ve her kenarı tam bölmeli: \(96=2^5\cdot 3\), \(120=2^3\cdot 3\cdot 5\) \(\Rightarrow\) \(\text{EBOB}(96,120)=2^3\cdot 3=24\) m.\\
Çevre: \(2\cdot(96+120)=432\) m. \ Kapalı çevre: ağaç sayısı \(=432\div 24=18\) (\(+1\) eklenmez).\\
Kenar kenar kontrol: \(96\div 24=4\), \(120\div 24=5\) aralık \(\Rightarrow\) \(2\cdot(4+5)=18\) aralık \(=18\) ağaç.
\end{cozum}

\begin{alistirma}
Üç otobüs hattından otobüsler sırasıyla \(20\), \(30\) ve \(45\) dakikada bir kalkmaktadır. Üç hattan saat 07{:}15'te birlikte otobüs kalktıktan sonra en erken saat kaçta yeniden birlikte kalkar?
\end{alistirma}
\begin{cozum}
Birlikte kalkış aralığı üç sürenin en küçük ortak katıdır:
\(20=2^2\cdot 5\), \(30=2\cdot 3\cdot 5\), \(45=3^2\cdot 5\) \(\Rightarrow\) \(\text{EKOK}=2^2\cdot 3^2\cdot 5=180\) dakika \(=3\) saat; \ 07{:}15\(\,+\,3\) saat \(=\) 10{:}15.
\end{cozum}

\begin{alistirma}
\(4\), \(6\) ve \(9\) ile bölündüğünde sırasıyla \(2\), \(4\) ve \(7\) kalanını veren üç basamaklı en küçük doğal sayı kaçtır?
\end{alistirma}
\begin{cozum}
Her kalan bölenden \(2\) eksiktir: \(4-2=2\), \(6-2=4\), \(9-2=7\). O halde \(x+2\) sayısı \(4\), \(6\) ve \(9\)'un ortak katıdır.\\
\(\text{EKOK}(4,6,9)=2^2\cdot 3^2=36\) \(\Rightarrow\) \(x=36k-2\): \ \(34,\ 70,\ 106,\dots\)\\
Üç basamaklı en küçük: \(x=106\). \ Kontrol: \(106=4\cdot 26+2=6\cdot 17+4=9\cdot 11+7\).
\end{cozum}

\kendinizoncesi
\section{Alıştırmalar (Kendiniz Çözün)}

\begin{enumerate}[label=\textbf{\arabic*.},leftmargin=*,itemsep=4pt]
\item \(\ab{24a}\) üç basamaklı sayısı \(4\) ile bölünebiliyorsa \(a\) hangi değerleri alabilir?
\item \(\ab{5a8}\) üç basamaklı sayısı \(3\) ile bölünebiliyorsa \(a\)'nın alabileceği değerlerin toplamı kaçtır?
\item \(9\,876\,543\) sayısının \(9\) ile ve \(5\) ile bölümünden kalanlar kaçtır?
\item \(3\,125\,764\) sayısının \(8\) ile bölümünden kalan kaçtır?
\item \(\text{EBOB}(72,108)\) ve \(\text{EKOK}(72,108)\) kaçtır?
\item İki pozitif tam sayının EBOB'u \(9\), EKOK'u \(270\)'tir. Sayılardan biri \(54\) ise diğeri kaçtır?
\item \(\ab{a4b}\) biçimindeki üç basamaklı sayılardan \(6\) ile bölünebilen en büyüğü kaçtır?
\item Boyutları \(150\) cm \(\times\) \(210\) cm olan dikdörtgen bir zemin, eş kare fayanslarla boşluksuz ve fayans kesilmeden döşenecektir. En az kaç fayans gerekir?
\item Kenarları \(15\) cm ve \(20\) cm olan eş dikdörtgen kartonlar, aynı yönde ve boşluksuz yan yana dizilerek bir kare oluşturulacaktır. En az kaç karton gerekir?
\item \(\ab{a1b}\) biçimindeki, \(45\) ile bölünebilen üç basamaklı sayıların toplamı kaçtır?
\item \(a<b\) olmak üzere EBOB'u \(4\), EKOK'u \(60\) olan kaç farklı \((a,b)\) pozitif tam sayı ikilisi vardır?
\item Bir kutudaki kalemler \(8\)'er, \(12\)'şer ve \(15\)'er sayıldığında her seferinde \(3\) kalem artmaktadır. Kalem sayısı \(200\)'den az olduğuna göre kutuda en çok kaç kalem vardır?
\end{enumerate}

\subsection*{Cevap Anahtarı}
{\small
\begin{multicols}{2}
\begin{enumerate}[label=\textbf{\arabic*.},leftmargin=*,itemsep=1pt]
\item \(\ab{4a}\in\{40,44,48\}\Rightarrow a\in\{0,4,8\}\).
\item \(13+a\) \(3\)'ün katı \(\Rightarrow a\in\{2,5,8\}\); toplam \(15\).
\item Rakam toplamı \(42\): \(9\) ile kalan \(6\); \(5\) ile kalan \(3\).
\item Son üç basamak \(764=8\cdot 95+4\Rightarrow\) kalan \(4\).
\item \(72=2^3\cdot 3^2\), \(108=2^2\cdot 3^3\); EBOB \(36\), EKOK \(216\).
\item \(54\cdot b=9\cdot 270\Rightarrow b=45\).
\item \(a=9\); \(b\) çift, \(13+b\) \(3\)'ün katı; en büyük \(948\).
\item \(\text{EBOB}(150,210)=30\); \(5\cdot 7=35\) fayans.
\item \(\text{EKOK}(15,20)=60\); \(4\cdot 3=12\) karton.
\item \(b=0\Rightarrow 810\); \(b=5\Rightarrow 315\); toplam \(1125\).
\item \(a=4x\), \(b=4y\), \(xy=15\): \((4,60)\), \((12,20)\); \(2\) ikili.
\item \(x-3\), \(120\)'nin katı: \(3,\ 123\); en çok \(123\).
\end{enumerate}
\end{multicols}}


\vspace{6pt}
\noindent\textcolor{anamavi}{\rule{\textwidth}{0.5pt}}\\
{\small Gelecek hafta: Rasyonel sayılar ve ondalık sayılar (dört işlem, sıralama).}

\end{document}
